\documentclass[11pt,a4paper]{ctexart}
\usepackage[margin=2.2cm]{geometry}
\usepackage{amsmath,amssymb,amsthm,mathtools,booktabs,hyperref,microtype}
\hypersetup{colorlinks=true,linkcolor=black,urlcolor=blue}
\newtheorem{theorem}{定理}
\title{\bfseries 有效双峰间隔、第二极值与单异常点剔除定律}
\author{研究札记 · F3-PEAK}
\date{2026 年 9 月 29 日}
\begin{document}
\maketitle
\begin{abstract}
固定普通整数参数的双根 \(F_3\) 轨道中，最高峰的精细上界需要深的 \(S\)-部分定理；本文证明第二高峰已经具有完全初等、全范围有效的 \(\log_3N+O(1)\) 界。核心是归一化二次多项式两点差分的因式分解与根的负有理迹。由此得到同门槛两峰的显式分离、第二次命中的有效指数尺度、一个峰对另一根类的反射定位、完整 \(3^T\) 区段的双剔除恒等式，以及负整数迹时删除一个最大值后的完全固定直方图。另得超临界峰的有效 \(O(\log\log X)\) 稀疏界。状态为 \texttt{PAPER-AUDITED}，尚未 Lean 形式化。
\end{abstract}

\section{设置与差分恒等式}
固定
\[
q=3^k,\qquad k\ge1,
\]
以及普通整数 \(n,d>1\)，满足基础双根条件。令
\[
Q_j=(n+2jqd)(n+(2jq+2)d)+1,
\]
\[
D_j=v_3(Q_j),\qquad e_j=D_j-2k.
\]
归一化
\[
G(X)=Q_X/q^2
=
4d^2X^2+4dKX+Q_0/q^2,
\qquad K=(n+d)/q.
\]
写
\[
g=\gcd(K,d),\qquad
A=K/g,\quad B=d/g.
\]
则两根迹为
\[
r_0+r_1=-A/B,
\qquad A,B>0,\quad3\nmid B.
\]

直接相减得到
\[
\boxed{
G(j)-G(i)
=
4dg(j-i)(B(i+j)+A).
}
\]
前因子是三进单位。

\section{双峰有效分离}
取
\[
0\le i<j,\qquad
t=\min(e_i,e_j).
\]
当 \(t\ge1\) 时，\(i,j\) 必在模 \(3\) 的两个根类中。

若同支：
\[
\boxed{3^t\mid j-i.}
\]
若异支：
\[
\boxed{3^t\mid B(i+j)+A.}
\]
若 \(e_i\ne e_j\)，相应整数的赋值恰为 \(t\)。

由于 \(A,B>0\) 且索引非负，
\[
0<j-i\le B(i+j)+A.
\]
所以统一有
\[
\boxed{
3^{\min(e_i,e_j)}
\le
A+B(i+j).
}
\]

\section{第二极值}
对 \(N\ge3\)，记前 \(N\) 项最大、第二大深度为
\[
M_1(N),\qquad M_2(N).
\]
令
\[
t=\lfloor\log_3N\rfloor.
\]
两个根模 \(3^t\) 的最小非负代表都小于 \(3^t\le N\)，且互异，所以
\[
M_2(N)\ge2k+t.
\]
另一方面，取达到两个最大值的不同索引 \(i<j<N\)，由双峰界
\[
3^{M_2(N)-2k}
\le
A+B(2N-3).
\]
因此
\[
\boxed{
2k+\lfloor\log_3N\rfloor
\le
M_2(N)
\le
2k+
\left\lfloor
\log_3(A+B(2N-3))
\right\rfloor.
}
\]
故
\[
M_2(N)=2k+\log_3N+O_{A,B}(1)
\]
且常数完全有效。

\section{第二次命中与根反射}
令 \(P=3^t\)，两根的最小代表为 \(a_0(t),a_1(t)\)。第二次达到门槛
\[
D_j\ge2k+t
\]
的位置是
\[
\tau_2(t)=\max(a_0(t),a_1(t)).
\]
由
\[
P\mid B(a_0+a_1)+A
\]
得到
\[
\boxed{
\max\left\{0,\frac{3^t-A}{2B}\right\}
\le
\tau_2(t)
<
3^t.
}
\]

若已知一个峰 \(a\) 满足 \(e_a\ge t\)，则另一根类由 Vieta 迹直接给出：
\[
\boxed{
b\equiv-a-AB^{-1}\pmod{3^t}.
}
\]
并有精确集合等式
\[
\boxed{
\{j\ge0:e_j\ge t\}
=
\{j\equiv a\pmod{3^t}\}
\sqcup
\{j\equiv b\pmod{3^t}\}.
}
\]

\section{完整 \(3^T\) 区段}
令
\[
P=3^T.
\]
对每个 \(1\le t\le T\)，区段 \(0\le j<P\) 中恰有
\[
2P/3^t
\]
个索引满足 \(e_j\ge t\)。因此
\[
\sum_{j<P}\min(e_j,T)=P-1.
\]
恰有两个索引达到 \(e_j\ge T\)，故
\[
\boxed{
\sum_{j<P}D_j-M_1(P)-M_2(P)
=
(2k+1)P-1-4k-2T.
}
\]

若根迹为负整数，即 \(B=1\)，且
\[
3^T>A-3,
\]
则第二峰恰为
\[
\boxed{M_2(3^T)=2k+T.}
\]
删除一个最大值后，剩余直方图完全固定：
\[
\begin{array}{c|c}
D&\text{次数}\\ \hline
2k&3^{T-1}\\
2k+t,\ 1\le t<T&4\cdot3^{T-t-1}\\
2k+T&1.
\end{array}
\]
所以
\[
\boxed{
\sum_{j<3^T}D_j-M_1(3^T)
=
(2k+1)(3^T-1)-T.
}
\]

\section{任意高孤峰与超临界稀疏性}
取
\[
n=(2q-1)d.
\]
此时迹恒为 \(-2\)。多项式
\[
(4q^2-1)d^2+1
\]
在 \(d\equiv1\pmod3\) 有简单根。Hensel 提升并在下一位选择不继续提升的分支，可令
\[
D_0=H
\]
达到任意预先指定的 \(H>2k+T\)，而其余前 \(3^T\) 项最大值仍为 \(2k+T\)。

固定 \(\theta>1\)。若
\[
a<b,\qquad
3^{e_a}>a^\theta,\quad3^{e_b}>b^\theta,
\]
则
\[
\boxed{
b>
\frac{a^\theta-A}{B}-a.
}
\]
取
\[
\beta=\frac{\theta+1}{2}
\]
和显式 \(J_0\)，可推出相邻超临界峰最终满足
\[
b>a^\beta.
\]
于是
\[
\boxed{
\#\{J_0\le j<X:3^{e_j}>j^\theta\}
\le
1+
\left\lfloor
\frac{\log(\log X/\log J_0)}
{\log\beta}
\right\rfloor.
}
\]

\section{边界}
有效 \(O(\log\log X)\) 稀疏性不能推出超临界集合有限；最高峰的最终有限性仍依赖更深的 \(S\)-部分/Roth 型输入。本文没有证明最高峰的 \(\log N+O(1)\) 界，也没有固定间距素数结论。

\begin{thebibliography}{9}
\bibitem{BEG} Y. Bugeaud, J.-H. Evertse and K. Győry, \emph{S-parts of values of univariate polynomials, binary forms and decomposable forms at integral points}, Acta Arith. 184 (2018).
\bibitem{Hensel} K. Conrad, \emph{Hensel's Lemma}.
\end{thebibliography}
\end{document}
